\documentclass[../DoAn.tex]{subfiles}
\begin{document}

Chương này trình bày toàn bộ quá trình em phân tích thiết kế, triển khai và đánh giá hệ thống NihongoFlow, dựa trên các yêu cầu đã xác định ở Chương \ref{chapter:Related_works} và các công nghệ đã lựa chọn ở Chương \ref{chapter:Methodology}.

\section{Thiết kế kiến trúc}

\subsection{Lựa chọn kiến trúc phần mềm}
\label{subsection:4.1.1}

Em thiết kế hệ thống theo kiến trúc ba lớp kết hợp với mô hình phân tách frontend--backend qua RESTful API. Phía backend, em tổ chức ba lớp như sau:

\begin{itemize}
    \item \textbf{Lớp Controller} (gói \texttt{controller}): tiếp nhận request HTTP, thực hiện validate đầu vào (\texttt{@Valid}) và ánh xạ sang DTO; không chứa logic nghiệp vụ.
    \item \textbf{Lớp Service} (gói \texttt{service}): chứa toàn bộ logic nghiệp vụ -- tính điểm quiz, kiểm tra điều kiện đăng ký khóa học, cập nhật streak, xác thực chữ ký VNPay, v.v.
    \item \textbf{Lớp Repository} (gói \texttt{repository}): các interface kế thừa \texttt{JpaRepository}, chỉ chứa định nghĩa truy vấn (kể cả truy vấn tùy biến bằng \texttt{@Query}), không chứa logic.
\end{itemize}

Bên cạnh ba lớp chính, em bổ sung các gói hỗ trợ: \texttt{entity} (các thực thể JPA ánh xạ tới bảng CSDL), \texttt{dto} (đối tượng truyền dữ liệu cho request/response), \texttt{security} (bộ lọc JWT, cấu hình \texttt{UserDetailsService}), \texttt{config} (cấu hình Spring, CORS, VNPay), và \texttt{exception} (xử lý lỗi tập trung qua \texttt{@RestControllerAdvice} và \texttt{ApiException}).

Phía frontend, em tổ chức kiến trúc tương ứng theo các tầng: Page (hiển thị giao diện, điều hướng), Hook (gói \texttt{hooks} -- toàn bộ \texttt{useQuery}/\texttt{useMutation} bằng TanStack Query), Lib (gói \texttt{lib} -- hàm thuần, hằng số, instance Axios), và Component (gói \texttt{components} -- thành phần giao diện tái sử dụng). Em quy định Page chỉ được phép gọi dữ liệu thông qua hook, không gọi trực tiếp Axios; hook không chứa state giao diện (như trạng thái mở/đóng modal).

So với mô hình MVC truyền thống, kiến trúc trên có thể coi lớp Controller + Service như phần "Controller" mở rộng, lớp Entity/Repository như "Model", và toàn bộ ứng dụng React như "View" -- giao tiếp với "Model/Controller" hoàn toàn qua API JSON thay vì render phía server.

\subsection{Thiết kế tổng quan}
\label{subsection:4.1.2}

Ở mức tổng quan nhất, em chia hệ thống thành ba gói lớn: \texttt{frontend}, \texttt{backend} và cơ sở dữ liệu, minh họa trong Hình \ref{fig:package-diagram-tongquan}.

\begin{figure}[H]
    \centering
    \includegraphics[width=0.7\textwidth]{Hinhve/package-diagram-tongquan.png}
    \caption{Biểu đồ gói tổng quan hệ thống NihongoFlow}
    \label{fig:package-diagram-tongquan}
\end{figure}

\textbf{Gói \texttt{frontend}.} \textit{Mục đích:} là gói chịu trách nhiệm hiển thị giao diện và tiếp nhận tương tác của cả hai nhóm người dùng (học viên, quản trị viên) -- cửa ngõ duy nhất mà người dùng tiếp xúc trực tiếp. \textit{Nhiệm vụ:}
\begin{itemize}
    \item Hiển thị nội dung học (khóa học, bài học, quiz) và khu vực quản trị dưới dạng giao diện web.
    \item Gửi yêu cầu (API call) tới \texttt{backend} cho mọi thao tác cần dữ liệu hoặc thay đổi trạng thái: đăng nhập, đăng ký khóa học, nộp bài quiz, quản trị nội dung, v.v.
    \item Giao tiếp một chiều với \texttt{backend} qua RESTful API -- không tự xử lý nghiệp vụ, không truy cập trực tiếp cơ sở dữ liệu.
\end{itemize}

\textbf{Gói \texttt{backend}.} \textit{Mục đích:} là gói trung tâm xử lý toàn bộ logic nghiệp vụ, đóng vai trò cầu nối giữa \texttt{frontend} và cơ sở dữ liệu, đồng thời là nơi duy nhất giao tiếp với các dịch vụ bên thứ ba. \textit{Nhiệm vụ:}
\begin{itemize}
    \item Xử lý các luồng nghiệp vụ: xác thực JWT, phân quyền RBAC, chấm điểm quiz, tính streak, xử lý thanh toán.
    \item Giao tiếp với cơ sở dữ liệu để đọc/ghi dữ liệu thông qua Spring Data JPA/Hibernate.
    \item Gọi tới cổng thanh toán VNPay để sinh URL thanh toán và xác thực callback; phía \texttt{frontend} gọi trực tiếp YouTube IFrame API để phát video (không đi qua \texttt{backend}).
    \item Đảm bảo \texttt{frontend} không cần biết cấu trúc dữ liệu nội bộ hay quy tắc nghiệp vụ.
\end{itemize}

\textbf{Gói cơ sở dữ liệu (MySQL).} \textit{Mục đích:} lưu trữ bền vững toàn bộ dữ liệu của hệ thống, đảm bảo toàn vẹn tham chiếu giữa các thực thể có quan hệ phân cấp chặt chẽ (Cấp độ -- Khóa học -- Bài học -- Quiz). \textit{Nhiệm vụ:}
\begin{itemize}
    \item Lưu trữ thông tin người dùng, nội dung học, lượt làm quiz, tiến độ học tập, giao dịch thanh toán.
    \item Đảm bảo tính nhất quán qua ràng buộc khóa ngoại và \texttt{ON DELETE CASCADE}.
    \item Chỉ được truy cập gián tiếp thông qua \texttt{backend} (qua JPA), không giao tiếp trực tiếp với \texttt{frontend}.
\end{itemize}

\textbf{Mối quan hệ giữa các gói.}
\begin{itemize}
    \item \texttt{frontend} $\to$ \texttt{backend}: mọi chức năng của người dùng (đăng nhập, đăng ký khóa học, làm quiz, quản trị nội dung...) đều được gửi về \texttt{backend} xử lý và trả kết quả qua JSON.
    \item \texttt{backend} $\to$ cơ sở dữ liệu: \texttt{backend} phụ thuộc cơ sở dữ liệu để lưu trữ và truy xuất toàn bộ dữ liệu nghiệp vụ, thực hiện thông qua Spring Data JPA/Hibernate.
    \item \texttt{backend} $\to$ VNPay: phụ thuộc một chiều ra dịch vụ bên thứ ba, chỉ để sinh URL thanh toán và xác thực chữ ký callback -- không có chiều phụ thuộc ngược lại.
\end{itemize}

Đi sâu hơn vào cấu trúc bên trong hai gói \texttt{frontend} và \texttt{backend}, Hình \ref{fig:package-diagram} mô tả biểu đồ gói chi tiết, phân rã mỗi gói lớn thành các gói con theo đúng cấu trúc thư mục mã nguồn và kiến trúc phân tầng ba lớp đã trình bày ở Mục \ref{subsection:4.1.1}. Phía frontend gồm \texttt{routes} (định tuyến bằng React Router trong \texttt{App.tsx}) dẫn tới \texttt{pages} -- chia thành ba nhóm theo vai trò người dùng: \textit{học viên}, \textit{quản trị}, \textit{xác thực} -- cùng \texttt{hooks} (tầng gọi API, tương ứng vai trò của gói \texttt{services} trong các framework khác) và \texttt{components}. Phía backend -- nơi tập trung phần lớn logic nghiệp vụ -- em vẽ chi tiết tới từng lớp con của hai gói \texttt{controller} (đủ 10 lớp Controller) và \texttt{service} (đủ 14 lớp Service), cùng các gói hỗ trợ \texttt{security} (bộ lọc JWT, đóng vai trò tương đương \texttt{middlewares}), \texttt{config}, \texttt{repository}, \texttt{entity}, \texttt{dto}, \texttt{exception}. Hệ thống không có phân hệ học máy hay tác vụ định kỳ (cron job) như một số đồ án tham khảo khác, nên biểu đồ không có các gói tương ứng.

\begin{figure}[H]
    \centering
    \includegraphics[width=\textwidth]{Hinhve/package-diagram.png}
    \caption{Biểu đồ gói chi tiết hệ thống NihongoFlow}
    \label{fig:package-diagram}
\end{figure}

\subsubsection{Gói backend}

\begin{itemize}
    \item \texttt{controller}: tiếp nhận request HTTP, validate đầu vào bằng \texttt{@Valid}, ánh xạ sang DTO và gọi \texttt{service} tương ứng -- không chứa logic nghiệp vụ.
    \item \texttt{service}: tầng nghiệp vụ chính -- chấm điểm quiz, kiểm tra điều kiện đăng ký khóa học, cập nhật streak, xác thực chữ ký VNPay, v.v.
    \item \texttt{repository}: các interface kế thừa \texttt{JpaRepository}, chỉ định nghĩa truy vấn (kể cả truy vấn tùy biến bằng \texttt{@Query}), không chứa logic.
    \item \texttt{entity}: các lớp JPA ánh xạ trực tiếp tới bảng cơ sở dữ liệu (\texttt{Question}, \texttt{QuizAttempt}, \texttt{User}, v.v.).
    \item \texttt{dto}: đối tượng truyền dữ liệu cho request/response, đảm bảo entity không bao giờ lộ trực tiếp qua API.
    \item \texttt{security}: bộ lọc xác thực JWT (\texttt{JwtAuthenticationFilter}), sinh/giải mã token (\texttt{JwtService}), cấu hình \texttt{UserDetailsService}.
    \item \texttt{config}: cấu hình Spring (\texttt{SecurityConfig}, CORS, thuộc tính VNPay, cookie).
    \item \texttt{exception}: xử lý lỗi tập trung qua \texttt{@RestControllerAdvice} (\texttt{GlobalExceptionHandler}) và \texttt{ApiException}, đảm bảo mọi lỗi trả về đúng envelope chuẩn.
\end{itemize}

\subsubsection{Gói frontend}

\begin{itemize}
    \item \texttt{pages}: hiển thị giao diện và xử lý điều hướng cho từng trang -- chỉ lấy dữ liệu thông qua \texttt{hooks}, không gọi Axios trực tiếp.
    \item \texttt{hooks}: toàn bộ \texttt{useQuery}/\texttt{useMutation} (TanStack Query) tổ chức theo domain (\texttt{useCourse}, \texttt{useQuiz}, \texttt{useDashboard}, v.v.), xử lý cache và invalidation.
    \item \texttt{lib}: hàm thuần, hằng số dùng chung và instance Axios duy nhất (\texttt{axios.ts}).
    \item \texttt{components}: thành phần giao diện tái sử dụng (\texttt{VideoPlayer}, các bước của quiz wizard, \texttt{Navbar}, v.v.).
    \item \texttt{store}: trạng thái xác thực phía client quản lý bằng Zustand (\texttt{authStore.ts}).
    \item \texttt{types}: toàn bộ khai báo kiểu TypeScript dùng chung giữa các gói trên.
\end{itemize}

Em thiết kế quan hệ phụ thuộc giữa các gói con tuân theo nguyên tắc một chiều: \texttt{controller} phụ thuộc \texttt{service}, \texttt{service} phụ thuộc \texttt{repository} và \texttt{entity}; không có chiều phụ thuộc ngược lại. Tương tự, \texttt{pages} phụ thuộc \texttt{hooks} và \texttt{components}, \texttt{hooks} phụ thuộc \texttt{lib} và \texttt{store}, không có gói tầng dưới phụ thuộc ngược lên tầng trên.

Hai chiều giao tiếp chính giữa frontend và backend là: (i) các request CRUD/nghiệp vụ qua \texttt{lib/axios.ts} (instance Axios duy nhất, có interceptor tự động đính kèm access token và làm mới token khi hết hạn); (ii) request callback từ VNPay tới \texttt{PaymentController}, không đi qua frontend.

\subsection{Thiết kế chi tiết gói}
\label{subsection:4.1.3}

Hai nhóm package em đánh giá có quan hệ phức tạp nhất trong hệ thống là nhóm Quiz và nhóm Enrollment, được em trình bày chi tiết dưới đây.

\textbf{Nhóm Quiz.} Hình \ref{fig:class-diagram-quiz} mô tả các lớp \texttt{Quiz}, \texttt{Question}, \texttt{QuizAttempt} và quan hệ giữa chúng. Một \texttt{Lesson} có quan hệ 1--1 với \texttt{Quiz}; một \texttt{Quiz} kết hợp (composition) với nhiều \texttt{Question} (đúng 10 câu, chia theo \texttt{questionType}); mỗi \texttt{QuizAttempt} liên kết tới một \texttt{Quiz} và một \texttt{User}, lưu kết quả dưới dạng cột JSON \texttt{answers}. \texttt{QuizService} phụ thuộc (dependency) vào cả ba repository tương ứng để thực hiện việc chấm điểm.

\begin{figure}[H]
    \centering
    \includegraphics[width=0.85\textwidth]{Hinhve/class-diagram-quiz.png}
    \caption{Biểu đồ lớp nhóm Quiz}
    \label{fig:class-diagram-quiz}
\end{figure}

\textbf{Nhóm Enrollment.} Hình \ref{fig:class-diagram-enrollment} mô tả quan hệ giữa \texttt{User}, \texttt{Course} và \texttt{UserCourseEnrollment} -- thực thể trung gian biểu diễn quan hệ nhiều--nhiều giữa người dùng và khóa học, có khóa chính phức hợp (\texttt{userId}, \texttt{courseId}). \texttt{Payment} liên kết với \texttt{User} và \texttt{Course}, được \texttt{PaymentService} sử dụng để tạo và xác thực giao dịch VNPay; khi giao dịch thành công, \texttt{PaymentService} gọi \texttt{EnrollmentService} (kết hợp/association) để tạo bản ghi \texttt{UserCourseEnrollment}.

\begin{figure}[H]
    \centering
    \includegraphics[width=0.85\textwidth]{Hinhve/class-diagram-enrollment.png}
    \caption{Biểu đồ lớp nhóm Enrollment và Payment}
    \label{fig:class-diagram-enrollment}
\end{figure}

\section{Thiết kế chi tiết}

\subsection{Thiết kế giao diện}
\label{subsection:4.2.1}

Em thiết kế giao diện hệ thống hướng tới máy tính để bàn, độ phân giải tham chiếu \(1920\times1080\), với bố cục nội dung giới hạn chiều rộng tối đa \(1120\)px, căn giữa. Em lựa chọn phong cách thiết kế tối giản, nền tối, lấy cảm hứng từ Resend.com, với bộ màu chuẩn hóa qua biến CSS:

\begin{itemize}
    \item \textbf{Nền}: \texttt{--bg} (\#0a0a0a), \texttt{--s1} (\#111111) cho thẻ/khối nội dung, \texttt{--s2} (\#161616) cho khối lồng bên trong
    \item \textbf{Viền}: \texttt{--b1} (\#262626), \texttt{--b2} (\#333333)
    \item \textbf{Chữ}: \texttt{--t1} (\#ededed) chữ chính, \texttt{--t2} (\#a1a1a1) chữ phụ, \texttt{--t3} (\#555555) chữ mờ
    \item \textbf{Nhấn}: \texttt{--acc} (\#e94560) -- duy nhất một màu nhấn cho toàn hệ thống, dùng cho nút hành động chính, thanh tiến độ và trạng thái được chọn
\end{itemize}

Phông chữ chính em chọn là Inter cho văn bản tiếng Việt/tiếng Anh và Noto Sans JP cho văn bản tiếng Nhật (cỡ chữ tối thiểu 14px, độ đậm tối thiểu 500 để đảm bảo hiển thị rõ nét). Khoảng cách giữa các phần tử tuân theo lưới 8px. Các thành phần dùng chung (nút, thẻ, huy hiệu, ô nhập liệu) em xây dựng từ shadcn/ui và tùy biến lại theo bộ token trên thông qua \texttt{globals.css}, đảm bảo không có giá trị màu hay khoảng cách "hard-code" rải rác trong mã nguồn.

Một số màn hình tiêu biểu minh họa các nguyên tắc trên gồm trang chủ (Hình \ref{fig:screenshot-landing}), trang khóa học, trang xem video kết hợp quiz, và khu vực quản trị, được em trình bày chi tiết tại Mục \ref{subsection:4.3.3}.

\subsection{Thiết kế lớp}
\label{subsection:4.2.2}

Phần này em trình bày chi tiết ba lớp chủ đạo của hệ thống: \texttt{Question}, \texttt{QuizAttempt}, và \texttt{UserCourseEnrollment}.

\textbf{Lớp \texttt{Question}} đại diện cho một câu hỏi trong quiz, với các thuộc tính: \texttt{id} (Long), \texttt{quiz} (quan hệ \texttt{ManyToOne} tới \texttt{Quiz}), \texttt{content} (String -- nội dung câu hỏi), \texttt{options} (kiểu JSON, mảng 4 chuỗi đáp án), \texttt{correctOption} (int -- chỉ số đáp án đúng từ 0--3, chỉ áp dụng cho loại CONTENT), \texttt{correctAnswerText} (String -- đáp án văn bản đúng, chỉ áp dụng cho loại VOCABULARY), \texttt{correctOrder} (kiểu JSON -- thứ tự đúng, chỉ áp dụng cho loại SEQUENCE), \texttt{orderIndex} (int -- thứ tự câu hỏi từ 1--10), và \texttt{questionType} (enum: \texttt{VOCABULARY}, \texttt{CONTENT}, \texttt{SEQUENCE}). Thiết kế tổng quát hóa qua \texttt{questionType} cho phép em thêm dạng câu hỏi mới trong tương lai mà không cần thay đổi schema -- nội dung này em phân tích sâu hơn ở Mục \ref{section:5.1}.

\textbf{Lớp \texttt{QuizAttempt}} lưu một lượt làm bài, gồm: \texttt{id}, \texttt{user} (\texttt{ManyToOne} tới \texttt{User}), \texttt{quiz} (\texttt{ManyToOne} tới \texttt{Quiz}), \texttt{score} (int), \texttt{passed} (boolean), \texttt{answers} (kiểu JSON -- ánh xạ \texttt{questionId} sang đáp án đã chọn, kiểu dữ liệu của đáp án phụ thuộc \texttt{questionType} của từng câu), và \texttt{attemptedAt} (thời điểm nộp bài). Phương thức nghiệp vụ chính \texttt{QuizService.submitAttempt(...)} nhận vào \texttt{quizId}, \texttt{answers} và người dùng hiện tại, thực hiện: (i) kiểm tra tập \texttt{questionId} gửi lên khớp chính xác (đối xứng hai chiều) với tập câu hỏi của quiz; (ii) so sánh từng đáp án với đáp án đúng tương ứng theo \texttt{questionType} để tính số câu đúng; (iii) so sánh với \texttt{Quiz.passScore} để xác định \texttt{passed}; (iv) lưu \texttt{QuizAttempt} và gọi \texttt{LessonService.markCompleted(...)} nếu \texttt{passed = true}, qua đó cập nhật streak nếu đây là lượt hoàn thành bài học đầu tiên trong ngày.

\textbf{Lớp \texttt{UserCourseEnrollment}} là thực thể có khóa chính phức hợp (\texttt{userId}, \texttt{courseId}), gồm thêm thuộc tính \texttt{enrolledAt}. Phương thức \texttt{PaymentService.createPayment(...)} kiểm tra điều kiện: nếu \texttt{Course.price == 0}, gọi trực tiếp \texttt{EnrollmentService.enroll(...)} để tạo bản ghi enrollment; nếu \texttt{price > 0}, việc enroll chỉ được thực hiện sau khi \texttt{Payment} được xác nhận thành công qua luồng VNPay.

Hai luồng tương tác em đánh giá là quan trọng nhất được minh họa bằng biểu đồ trình tự:

\begin{figure}[H]
    \centering
    \includegraphics[width=0.9\textwidth]{Hinhve/sequence-submit-quiz.png}
    \caption{Biểu đồ trình tự: Nộp bài quiz}
    \label{fig:sequence-submit-quiz}
\end{figure}

Hình \ref{fig:sequence-submit-quiz} mô tả luồng nộp bài quiz: \texttt{QuizPage} (frontend) gọi hook \texttt{useSubmitQuiz}, gửi \texttt{\seqsplit{POST /api/v1/quizzes/:id/attempts}} tới \texttt{QuizController}; \texttt{QuizController} gọi \texttt{\seqsplit{QuizService.submitAttempt}}, dịch vụ này truy vấn \texttt{QuestionRepository} để lấy đáp án đúng, tính điểm, lưu \texttt{QuizAttempt} qua \texttt{\seqsplit{QuizAttemptRepository}}, đồng thời gọi \texttt{LessonService}/\texttt{StreakService} để cập nhật trạng thái hoàn thành bài học và chuỗi ngày học; kết quả (điểm số, đạt/không đạt) được trả về và hiển thị qua popup kết quả.

\begin{figure}[H]
    \centering
    \includegraphics[width=0.9\textwidth]{Hinhve/sequence-enroll-payment.png}
    \caption{Biểu đồ trình tự: Đăng ký khóa học có phí qua VNPay}
    \label{fig:sequence-enroll-payment}
\end{figure}

Hình \ref{fig:sequence-enroll-payment} mô tả luồng đăng ký khóa học có phí: \texttt{CourseDetailPage} gọi \texttt{\seqsplit{useCreatePaymentMutation}}, gửi \texttt{\seqsplit{POST /api/v1/payments/create}} tới \texttt{PaymentController}; \texttt{PaymentService} tạo bản ghi \texttt{Payment} ở trạng thái chờ, gọi \texttt{VnpayService} để sinh URL thanh toán kèm chữ ký HMAC-SHA512, và trả URL về cho frontend để chuyển hướng (\texttt{\seqsplit{window.location.href}}). Sau khi người dùng thanh toán trên cổng VNPay, VNPay gọi lại \texttt{\seqsplit{GET /api/v1/payments/vnpay-return}}; \texttt{VnpayService} xác thực lại chữ ký, \texttt{PaymentService} cập nhật trạng thái \texttt{Payment} và gọi \texttt{\seqsplit{EnrollmentService.enroll(...)}} nếu thành công, sau đó \texttt{PaymentController} chuyển hướng (redirect) trình duyệt về trang \texttt{\seqsplit{/thanh-toan/ket-qua}} của frontend.

\subsection{Thiết kế cơ sở dữ liệu}
\label{subsection:4.2.3}

Hình \ref{fig:erd} trình bày biểu đồ thực thể liên kết (ERD) của toàn bộ hệ thống, gồm 12 bảng chính: \texttt{users}, \texttt{levels}, \texttt{courses}, \texttt{lessons}, \texttt{quizzes}, \texttt{questions}, \texttt{quiz\_attempts}, \texttt{\seqsplit{user\_lesson\_progress}}, \texttt{\seqsplit{user\_course\_enrollments}}, \texttt{user\_streaks}, \texttt{payments}, và bảng lưu refresh token.

\begin{figure}[H]
    \centering
    \includegraphics[width=0.95\textwidth]{Hinhve/erd.png}
    \caption{Biểu đồ thực thể liên kết (ERD) hệ thống NihongoFlow}
    \label{fig:erd}
\end{figure}

Các bảng chính trong cơ sở dữ liệu, tương ứng với các thực thể trên ERD, gồm:
\begin{itemize}
    \item \texttt{users}: tài khoản người dùng (học viên hoặc quản trị viên).
    \item \texttt{refresh\_tokens}: refresh token đang hiệu lực của từng người dùng, phục vụ cơ chế xoay vòng token (Mục \ref{section:5.3}).
    \item \texttt{user\_streaks}: chuỗi ngày học liên tục hiện tại và dài nhất của từng người dùng.
    \item \texttt{levels}: cấp độ tiếng Nhật (N5--N1).
    \item \texttt{courses}: khóa học, thuộc một cấp độ.
    \item \texttt{lessons}: bài học (video), thuộc một khóa học.
    \item \texttt{quizzes}: bài kiểm tra sau mỗi bài học -- quan hệ 1--1 với \texttt{lessons}.
    \item \texttt{questions}: câu hỏi của một quiz, tổng quát hóa theo \texttt{question\_type} (Mục \ref{section:5.1}).
    \item \texttt{quiz\_attempts}: lượt làm bài quiz của học viên, lưu điểm số và đáp án đã chọn.
    \item \texttt{user\_lesson\_progress}: tiến độ xem video của từng học viên cho từng bài học.
    \item \texttt{user\_course\_enrollments}: đăng ký khóa học của học viên.
    \item \texttt{payments}: giao dịch thanh toán VNPay cho các khóa học có phí.
\end{itemize}

Một số điểm thiết kế em lưu ý đặc biệt khi xây dựng schema:

\begin{itemize}
    \item \textbf{Khóa chính phức hợp}: \texttt{\seqsplit{user\_lesson\_progress}} dùng khóa chính phức hợp (\texttt{user\_id}, \texttt{lesson\_id}); \texttt{\seqsplit{user\_course\_enrollments}} dùng khóa chính phức hợp (\texttt{user\_id}, \texttt{course\_id}). Thiết kế này đảm bảo mỗi người dùng chỉ có đúng một bản ghi tiến độ cho mỗi bài học và đúng một bản ghi đăng ký cho mỗi khóa học.
    \item \textbf{Cột kiểu JSON}: \texttt{questions.options} (mảng 4 chuỗi đáp án), \texttt{questions.correct\_order} (thứ tự đúng cho câu trình tự) và \texttt{\seqsplit{quiz\_attempts.answers}} (ánh xạ \texttt{questionId} sang đáp án đã chọn) được em lưu dưới dạng JSON, cho phép linh hoạt số lượng và cấu trúc mà không cần bảng phụ.
    \item \textbf{Phân cấp nội dung}: quan hệ một--nhiều theo chuỗi \texttt{levels} \(\to\) \texttt{courses} \(\to\) \texttt{lessons} \(\to\) \texttt{quizzes} \(\to\) \texttt{questions}, mỗi bảng con có khóa ngoại trỏ tới bảng cha và cột \texttt{order}/\texttt{order\_index} để xác định thứ tự hiển thị.
    \item \textbf{Ràng buộc xóa lan truyền (\texttt{ON DELETE CASCADE})}: khi một \texttt{lesson} hoặc \texttt{course} bị xóa, các bản ghi liên quan trong các bảng \texttt{quizzes}, \texttt{questions}, \texttt{quiz\_attempts}, \texttt{\seqsplit{user\_lesson\_progress}} và \texttt{\seqsplit{user\_course\_enrollments}} được tự động xóa theo, tránh dữ liệu mồ côi.
    \item \textbf{Quản lý phiên bản schema}: toàn bộ cấu trúc CSDL được em quản lý qua Flyway, cho phép áp dụng schema một cách có kiểm soát và tái lập trên nhiều môi trường. So với một schema quan hệ tối thiểu, lược đồ của hệ thống được em mở rộng thêm: bảng \texttt{\seqsplit{user\_course\_enrollments}} (đăng ký khóa học), cờ \texttt{hidden} và cột \texttt{price} trên \texttt{courses} (ẩn/hiện và định giá khóa học), bảng \texttt{payments} (lịch sử thanh toán), cột \texttt{question\_type} cùng hai cột đáp án riêng \texttt{correct\_answer\_text}/\texttt{correct\_order} trên \texttt{questions} (hỗ trợ ba dạng câu hỏi VOCABULARY/CONTENT/SEQUENCE với cách chấm điểm khác nhau -- xem Mục \ref{section:5.1}), và ràng buộc \texttt{ON DELETE CASCADE} đã nêu ở trên.
\end{itemize}

Bảng \ref{table:db-users} đến Bảng \ref{table:db-payments} mô tả chi tiết thuộc tính của từng bảng trong cơ sở dữ liệu.

\begin{table}[H]
\centering
\small
\renewcommand{\arraystretch}{1.2}
\begin{tabular}{|p{3.2cm}|p{2.4cm}|p{8cm}|}
\hline
\textbf{Thuộc tính} & \textbf{Kiểu dữ liệu} & \textbf{Mô tả} \\ \hline
id & BIGINT & Khóa chính (PK), tự tăng. \\ \hline
email & VARCHAR & Địa chỉ email, dùng để đăng nhập (duy nhất -- UNIQUE, không rỗng -- NN). \\ \hline
password & VARCHAR & Mật khẩu đã băm bằng BCrypt (không rỗng -- NN). \\ \hline
role & VARCHAR & Vai trò người dùng: \texttt{STUDENT} hoặc \texttt{ADMIN} (không rỗng -- NN). \\ \hline
created\_at & TIMESTAMP & Thời điểm tạo tài khoản. \\ \hline
\end{tabular}
\renewcommand{\arraystretch}{1}
\caption{Mô tả bảng \texttt{users}}
\label{table:db-users}
\end{table}

\begin{table}[H]
\centering
\small
\renewcommand{\arraystretch}{1.2}
\begin{tabular}{|p{3.2cm}|p{2.4cm}|p{8cm}|}
\hline
\textbf{Thuộc tính} & \textbf{Kiểu dữ liệu} & \textbf{Mô tả} \\ \hline
id & BIGINT & Khóa chính (PK), tự tăng. \\ \hline
user\_id & BIGINT & Khóa ngoại (FK tới \texttt{users.id}), không rỗng -- NN. \\ \hline
token & VARCHAR & Chuỗi UUID ngẫu nhiên đại diện cho refresh token (duy nhất -- UNIQUE). \\ \hline
expires\_at & TIMESTAMP & Thời điểm hết hạn (7 ngày kể từ lúc cấp). \\ \hline
revoked & BOOLEAN & Đánh dấu token đã bị thu hồi hay chưa. \\ \hline
\end{tabular}
\renewcommand{\arraystretch}{1}
\caption{Mô tả bảng \texttt{refresh\_tokens}}
\label{table:db-refresh-tokens}
\end{table}

\begin{table}[H]
\centering
\small
\renewcommand{\arraystretch}{1.2}
\begin{tabular}{|p{3.2cm}|p{2.4cm}|p{8cm}|}
\hline
\textbf{Thuộc tính} & \textbf{Kiểu dữ liệu} & \textbf{Mô tả} \\ \hline
user\_id & BIGINT & Khóa chính \& khóa ngoại (PK, FK tới \texttt{users.id}) -- quan hệ 1--1 với \texttt{users}. \\ \hline
current\_streak & INT & Số ngày học liên tục hiện tại. \\ \hline
longest\_streak & INT & Số ngày học liên tục dài nhất từng đạt được. \\ \hline
last\_activity\_at & TIMESTAMP & Thời điểm hoàn thành bài học gần nhất, dùng để tính streak còn hiệu lực hay không. \\ \hline
\end{tabular}
\renewcommand{\arraystretch}{1}
\caption{Mô tả bảng \texttt{user\_streaks}}
\label{table:db-user-streaks}
\end{table}

\begin{table}[H]
\centering
\small
\renewcommand{\arraystretch}{1.2}
\begin{tabular}{|p{3.2cm}|p{2.4cm}|p{8cm}|}
\hline
\textbf{Thuộc tính} & \textbf{Kiểu dữ liệu} & \textbf{Mô tả} \\ \hline
id & BIGINT & Khóa chính (PK), tự tăng. \\ \hline
name & VARCHAR & Tên cấp độ (N5, N4, N3, N2, N1), không rỗng -- NN. \\ \hline
description & VARCHAR & Mô tả ngắn về cấp độ. \\ \hline
order\_index & INT & Thứ tự hiển thị (1 = N5 ... 5 = N1). \\ \hline
\end{tabular}
\renewcommand{\arraystretch}{1}
\caption{Mô tả bảng \texttt{levels}}
\label{table:db-levels}
\end{table}

\begin{table}[H]
\centering
\small
\renewcommand{\arraystretch}{1.2}
\begin{tabular}{|p{3.2cm}|p{2.4cm}|p{8cm}|}
\hline
\textbf{Thuộc tính} & \textbf{Kiểu dữ liệu} & \textbf{Mô tả} \\ \hline
id & BIGINT & Khóa chính (PK), tự tăng. \\ \hline
level\_id & BIGINT & Khóa ngoại (FK tới \texttt{levels.id}), không rỗng -- NN. \\ \hline
title & VARCHAR & Tên khóa học, không rỗng -- NN. \\ \hline
description & TEXT & Mô tả chi tiết khóa học. \\ \hline
thumbnail\_url & VARCHAR & URL ảnh đại diện khóa học. \\ \hline
order\_index & INT & Thứ tự hiển thị trong danh sách. \\ \hline
hidden & BOOLEAN & Ẩn/hiện khỏi danh sách khám phá khóa học -- không ảnh hưởng học viên đã đăng ký. \\ \hline
price & BIGINT & Giá tiền (VNĐ); \texttt{0} nghĩa là khóa học miễn phí. \\ \hline
\end{tabular}
\renewcommand{\arraystretch}{1}
\caption{Mô tả bảng \texttt{courses}}
\label{table:db-courses}
\end{table}

\begin{table}[H]
\centering
\small
\renewcommand{\arraystretch}{1.2}
\begin{tabular}{|p{3.2cm}|p{2.4cm}|p{8cm}|}
\hline
\textbf{Thuộc tính} & \textbf{Kiểu dữ liệu} & \textbf{Mô tả} \\ \hline
id & BIGINT & Khóa chính (PK), tự tăng. \\ \hline
course\_id & BIGINT & Khóa ngoại (FK tới \texttt{courses.id}), không rỗng -- NN. \\ \hline
title & VARCHAR & Tên bài học, không rỗng -- NN. \\ \hline
video\_url & VARCHAR & URL video YouTube của bài giảng. \\ \hline
duration & INT & Thời lượng video (giây), tự động đọc qua YouTube IFrame API. \\ \hline
order\_index & INT & Thứ tự bài học trong khóa học. \\ \hline
\end{tabular}
\renewcommand{\arraystretch}{1}
\caption{Mô tả bảng \texttt{lessons}}
\label{table:db-lessons}
\end{table}

\begin{table}[H]
\centering
\small
\renewcommand{\arraystretch}{1.2}
\begin{tabular}{|p{3.2cm}|p{2.4cm}|p{8cm}|}
\hline
\textbf{Thuộc tính} & \textbf{Kiểu dữ liệu} & \textbf{Mô tả} \\ \hline
id & BIGINT & Khóa chính (PK), tự tăng. \\ \hline
lesson\_id & BIGINT & Khóa ngoại (FK tới \texttt{lessons.id}, duy nhất -- UNIQUE) -- đảm bảo quan hệ 1--1 với \texttt{lessons}. \\ \hline
pass\_score & INT & Ngưỡng điểm đạt (\%) -- 70/75/80 tùy cấp độ (Mục \ref{section:2.4}). \\ \hline
\end{tabular}
\renewcommand{\arraystretch}{1}
\caption{Mô tả bảng \texttt{quizzes}}
\label{table:db-quizzes}
\end{table}

\begin{table}[H]
\centering
\small
\renewcommand{\arraystretch}{1.2}
\begin{tabular}{|p{3.2cm}|p{2.4cm}|p{8cm}|}
\hline
\textbf{Thuộc tính} & \textbf{Kiểu dữ liệu} & \textbf{Mô tả} \\ \hline
id & BIGINT & Khóa chính (PK), tự tăng. \\ \hline
quiz\_id & BIGINT & Khóa ngoại (FK tới \texttt{quizzes.id}), không rỗng -- NN. \\ \hline
content & TEXT & Nội dung câu hỏi, không rỗng -- NN. \\ \hline
options & JSON & Mảng 4 chuỗi đáp án. \\ \hline
correct\_option & INT & Chỉ số đáp án đúng (0--3) -- chỉ áp dụng cho loại \texttt{CONTENT}. \\ \hline
correct\_answer\_text & VARCHAR & Đáp án văn bản đúng -- chỉ áp dụng cho loại \texttt{VOCABULARY}. \\ \hline
correct\_order & JSON & Thứ tự đúng (mảng chỉ số 0--3) -- chỉ áp dụng cho loại \texttt{SEQUENCE}. \\ \hline
order\_index & INT & Vị trí câu hỏi trong quiz (1--10). \\ \hline
question\_type & VARCHAR & Loại câu hỏi: \texttt{VOCABULARY} / \texttt{CONTENT} / \texttt{SEQUENCE}. \\ \hline
\end{tabular}
\renewcommand{\arraystretch}{1}
\caption{Mô tả bảng \texttt{questions}}
\label{table:db-questions}
\end{table}

\begin{table}[H]
\centering
\small
\renewcommand{\arraystretch}{1.2}
\begin{tabular}{|p{3.2cm}|p{2.4cm}|p{8cm}|}
\hline
\textbf{Thuộc tính} & \textbf{Kiểu dữ liệu} & \textbf{Mô tả} \\ \hline
id & BIGINT & Khóa chính (PK), tự tăng. \\ \hline
user\_id & BIGINT & Khóa ngoại (FK tới \texttt{users.id}), không rỗng -- NN. \\ \hline
quiz\_id & BIGINT & Khóa ngoại (FK tới \texttt{quizzes.id}), không rỗng -- NN. \\ \hline
score & INT & Điểm số phần trăm (0--100). \\ \hline
answers & JSON & Ánh xạ \texttt{questionId} sang đáp án đã chọn, kiểu dữ liệu tùy \texttt{question\_type}. \\ \hline
attempted\_at & TIMESTAMP & Thời điểm nộp bài. \\ \hline
passed & BOOLEAN & Đạt hay không đạt so với \texttt{pass\_score} của quiz. \\ \hline
\end{tabular}
\renewcommand{\arraystretch}{1}
\caption{Mô tả bảng \texttt{quiz\_attempts}}
\label{table:db-quiz-attempts}
\end{table}

\begin{table}[H]
\centering
\small
\renewcommand{\arraystretch}{1.2}
\begin{tabular}{|p{3.2cm}|p{2.4cm}|p{8cm}|}
\hline
\textbf{Thuộc tính} & \textbf{Kiểu dữ liệu} & \textbf{Mô tả} \\ \hline
user\_id & BIGINT & Khóa chính \& khóa ngoại (PK, FK tới \texttt{users.id}). \\ \hline
lesson\_id & BIGINT & Khóa chính \& khóa ngoại (PK, FK tới \texttt{lessons.id}) -- cùng \texttt{user\_id} tạo thành khóa chính phức hợp. \\ \hline
completed & BOOLEAN & Đã xem hết video (hoặc đạt quiz) hay chưa. \\ \hline
watched\_seconds & INT & Vị trí xem gần nhất (giây), dùng để phát tiếp (resume). \\ \hline
updated\_at & TIMESTAMP & Thời điểm cập nhật tiến độ gần nhất. \\ \hline
\end{tabular}
\renewcommand{\arraystretch}{1}
\caption{Mô tả bảng \texttt{user\_lesson\_progress}}
\label{table:db-user-lesson-progress}
\end{table}

\begin{table}[H]
\centering
\small
\renewcommand{\arraystretch}{1.2}
\begin{tabular}{|p{3.2cm}|p{2.4cm}|p{8cm}|}
\hline
\textbf{Thuộc tính} & \textbf{Kiểu dữ liệu} & \textbf{Mô tả} \\ \hline
user\_id & BIGINT & Khóa chính \& khóa ngoại (PK, FK tới \texttt{users.id}). \\ \hline
course\_id & BIGINT & Khóa chính \& khóa ngoại (PK, FK tới \texttt{courses.id}) -- cùng \texttt{user\_id} tạo thành khóa chính phức hợp. \\ \hline
enrolled\_at & TIMESTAMP & Thời điểm đăng ký khóa học. \\ \hline
\end{tabular}
\renewcommand{\arraystretch}{1}
\caption{Mô tả bảng \texttt{user\_course\_enrollments}}
\label{table:db-user-course-enrollments}
\end{table}

\begin{table}[H]
\centering
\small
\renewcommand{\arraystretch}{1.2}
\begin{tabular}{|p{3.2cm}|p{2.4cm}|p{8cm}|}
\hline
\textbf{Thuộc tính} & \textbf{Kiểu dữ liệu} & \textbf{Mô tả} \\ \hline
id & BIGINT & Khóa chính (PK), tự tăng. \\ \hline
user\_id & BIGINT & Khóa ngoại (FK tới \texttt{users.id}), không rỗng -- NN. \\ \hline
course\_id & BIGINT & Khóa ngoại (FK tới \texttt{courses.id}), không rỗng -- NN. \\ \hline
amount & BIGINT & Số tiền giao dịch (VNĐ). \\ \hline
vnp\_txn\_ref & VARCHAR & Mã giao dịch duy nhất (UNIQUE) do hệ thống sinh, gửi kèm URL thanh toán VNPay. \\ \hline
vnp\_txn\_id & VARCHAR & Mã giao dịch phía VNPay trả về sau khi thanh toán. \\ \hline
status & VARCHAR & Trạng thái: \texttt{PENDING} / \texttt{SUCCESS} / \texttt{FAILED}. \\ \hline
created\_at & TIMESTAMP & Thời điểm tạo giao dịch. \\ \hline
paid\_at & TIMESTAMP & Thời điểm xác nhận thanh toán thành công (có thể rỗng nếu chưa/không thành công). \\ \hline
\end{tabular}
\renewcommand{\arraystretch}{1}
\caption{Mô tả bảng \texttt{payments}}
\label{table:db-payments}
\end{table}

\section{Xây dựng ứng dụng}

\subsection{Thư viện và công cụ sử dụng}
\label{subsection:4.3.1}

Bảng \ref{table:4.1} liệt kê các công cụ và thư viện chính em sử dụng để xây dựng hệ thống.

\begin{table}[H]
\centering
\renewcommand{\arraystretch}{1.2}
\begin{tabular}{p{4.3cm} p{6.2cm} p{3cm}}
\hline
\textbf{Thành phần} & \textbf{Công cụ / Thư viện} & \textbf{Phiên bản} \\ \hline
Ngôn ngữ backend & Java & 21 \\ \hline
Framework backend & Spring Boot & 3.x \\ \hline
Bảo mật & Spring Security + JJWT & 3.x / 0.12.x \\ \hline
ORM & Spring Data JPA (Hibernate) & 6.x \\ \hline
CSDL & MySQL & 8.0 \\ \hline
Quản lý migration & Flyway & 10.x \\ \hline
Build tool backend & Maven & 3.9.x \\ \hline
Ngôn ngữ frontend & TypeScript (strict mode) & 5.x \\ \hline
Framework frontend & React + Vite & React 18 / Vite 5 \\ \hline
Quản lý server state & TanStack Query & 5.x \\ \hline
Quản lý client state & Zustand & 4.x \\ \hline
Định tuyến & React Router & 6.x \\ \hline
CSS Framework & Tailwind CSS + shadcn/ui & 4.x \\ \hline
HTTP client & Axios & 1.x \\ \hline
Phát video & YouTube IFrame Player API & -- \\ \hline
Thanh toán & VNPay Sandbox & -- \\ \hline
Container hóa môi trường dev & Docker Compose (MySQL) & -- \\ \hline
\end{tabular}
\renewcommand{\arraystretch}{1}
\caption{Danh sách công cụ và thư viện sử dụng}
\label{table:4.1}
\end{table}

\subsection{Kết quả đạt được}
\label{subsection:4.3.2}

Em đã xây dựng và chạy được đầy đủ hai luồng người dùng (học viên và quản trị viên) trên môi trường phát triển cục bộ. Bảng \ref{table:4.2} thống kê quy mô mã nguồn theo từng tầng tại thời điểm hoàn thiện đồ án.

\begin{table}[H]
\centering
\renewcommand{\arraystretch}{1.1}
\begin{tabular}{llc}
\hline
\textbf{Phía} & \textbf{Thành phần} & \textbf{Số lượng} \\ \hline
\multirow{8}{*}{Backend} & Controller & 10 \\
 & Service & 14 \\
 & Repository & 12 \\
 & Entity & 15 \\
 & DTO & 27 \\
 & Lớp Security & 6 \\
 & Lớp Config & 5 \\
 & Lớp Exception & 3 \\ \hline
\multirow{5}{*}{Frontend} & Page & 13 \\
 & Custom hook & 11 \\
 & Lib util & 5 \\
 & Component dùng chung & 11 \\
 & Component theo domain (quiz, dashboard) & 9 \\ \hline
\end{tabular}
\caption{Thống kê quy mô mã nguồn theo tầng}
\label{table:4.2}
\end{table}

Dữ liệu mẫu cho môi trường phát triển bao gồm 5 cấp độ (N5--N1), 10 khóa học (2 khóa/cấp độ), 24 bài học, 24 quiz và 240 câu hỏi (mỗi quiz đúng 10 câu, chia 4 câu VOCABULARY, 3 câu CONTENT, 3 câu SEQUENCE).

Về mặt chức năng, em đã hoàn thiện: xác thực JWT (đăng ký, đăng nhập, làm mới token, đăng xuất); quản lý nội dung theo cấu trúc Cấp độ -- Khóa học -- Bài học -- Quiz; đăng ký khóa học (miễn phí và có phí qua VNPay); phát video YouTube với theo dõi tiến độ và tự động chuyển sang quiz khi xem hết video lần đầu; chấm điểm và lưu lịch sử làm quiz kèm chức năng xem lại đáp án theo từng loại câu hỏi; theo dõi streak học tập theo ngày; và toàn bộ chức năng quản trị (CRUD khóa học/bài học, quản lý câu hỏi quiz).

\subsection{Minh họa các chức năng chính}
\label{subsection:4.3.3}

Phần này trình bày ảnh chụp màn hình của các chức năng chính trong hệ thống.

\begin{figure}[H]
    \centering
    \includegraphics[width=0.9\textwidth]{Hinhve/screenshot-landing.png}
    \caption{Trang chủ -- giới thiệu các cấp độ N5--N1}
    \label{fig:screenshot-landing}
\end{figure}

\begin{figure}[H]
    \centering
    \includegraphics[width=0.9\textwidth]{Hinhve/screenshot-courses.png}
    \caption{Trang danh sách khóa học, lọc theo cấp độ}
    \label{fig:screenshot-courses}
\end{figure}

\begin{figure}[H]
    \centering
    \includegraphics[width=0.9\textwidth]{Hinhve/screenshot-course-detail.png}
    \caption{Trang chi tiết khóa học -- danh sách bài học và nút đăng ký/mua khóa học}
    \label{fig:screenshot-course-detail}
\end{figure}

\begin{figure}[H]
    \centering
    \includegraphics[width=0.9\textwidth]{Hinhve/screenshot-lesson-video.png}
    \caption{Trang xem video bài học với trình phát YouTube IFrame API}
    \label{fig:screenshot-lesson-video}
\end{figure}

\begin{figure}[H]
    \centering
    \includegraphics[width=0.9\textwidth]{Hinhve/screenshot-quiz.png}
    \caption{Trang làm bài quiz theo wizard 3 bước, chia theo 3 nhóm câu hỏi}
    \label{fig:screenshot-quiz}
\end{figure}

\begin{figure}[H]
    \centering
    \includegraphics[width=0.9\textwidth]{Hinhve/screenshot-dashboard.png}
    \caption{Trang Dashboard -- tiến độ học tập, streak và lịch sử làm quiz}
    \label{fig:screenshot-dashboard}
\end{figure}

\begin{figure}[H]
    \centering
    \includegraphics[width=0.9\textwidth]{Hinhve/screenshot-admin-courses.png}
    \caption{Trang quản trị -- danh sách và quản lý khóa học}
    \label{fig:screenshot-admin-courses}
\end{figure}

\begin{figure}[H]
    \centering
    \includegraphics[width=0.9\textwidth]{Hinhve/screenshot-admin-quiz.png}
    \caption{Trang quản trị -- chỉnh sửa câu hỏi quiz}
    \label{fig:screenshot-admin-quiz}
\end{figure}

\begin{figure}[H]
    \centering
    \includegraphics[width=0.9\textwidth]{Hinhve/screenshot-payment-result.png}
    \caption{Trang kết quả thanh toán VNPay}
    \label{fig:screenshot-payment-result}
\end{figure}

\section{Kiểm thử}
\label{section:4.4}

Em kiểm thử hệ thống bằng phương pháp kiểm thử thủ công theo kịch bản, tập trung vào các luồng nghiệp vụ cần dữ liệu thực và môi trường tích hợp đầy đủ (cơ sở dữ liệu, cổng thanh toán sandbox) mà kiểm thử tự động khó mô phỏng chính xác.

Em kiểm thử thủ công trên hai chức năng quan trọng nhất: nộp bài quiz và đăng ký khóa học có phí qua VNPay. Bảng \ref{table:4.3} và Bảng \ref{table:4.4} trình bày các trường hợp kiểm thử tiêu biểu.

\begin{table}[H]
\centering
\small
\renewcommand{\arraystretch}{1.3}
\begin{tabular}{p{2cm}p{4.2cm}p{4.2cm}p{1.3cm}}
\hline
\textbf{Mã TC} & \textbf{Đầu vào} & \textbf{Kết quả mong đợi} & \textbf{Kết quả} \\ \hline
TC-Quiz-01 & Nộp đủ 10 câu, tất cả đáp án đúng & \texttt{score = 100}, \texttt{passed = true} & Đạt \\ \hline
TC-Quiz-02 & Nộp đủ 10 câu, số câu đúng dưới ngưỡng \texttt{passScore} & \texttt{passed = false}, hiển thị nút "Làm lại" & Đạt \\ \hline
TC-Quiz-03 & Nộp thiếu một \texttt{questionId} so với quiz & API trả lỗi \seqsplit{VALIDATION\_ERROR} & Đạt \\ \hline
TC-Quiz-04 & Nộp dư một \texttt{questionId} không thuộc quiz & API trả lỗi \seqsplit{VALIDATION\_ERROR} & Đạt \\ \hline
TC-Quiz-05 & Lượt làm đầu tiên trong ngày, kết quả đạt & \seqsplit{UserStreak.currentStreak} tăng thêm 1 & Đạt \\ \hline
\end{tabular}
\renewcommand{\arraystretch}{1}
\caption{Trường hợp kiểm thử chức năng nộp bài quiz}
\label{table:4.3}
\end{table}

\begin{table}[H]
\centering
\small
\renewcommand{\arraystretch}{1.3}
\begin{tabular}{p{2cm}p{4.2cm}p{4.2cm}p{1.3cm}}
\hline
\textbf{Mã TC} & \textbf{Đầu vào} & \textbf{Kết quả mong đợi} & \textbf{Kết quả} \\ \hline
TC-Pay-01 & Đăng ký khóa học \texttt{price = 0} & Tạo bản ghi \seqsplit{UserCourseEnrollment} ngay, không qua VNPay & Đạt \\ \hline
TC-Pay-02 & Đăng ký khóa học \texttt{price > 0}, thanh toán thành công trên sandbox & Callback xác thực HMAC-SHA512 hợp lệ, tạo \texttt{Payment} thành công và \seqsplit{UserCourseEnrollment} & Đạt \\ \hline
TC-Pay-03 & Hủy giao dịch trên trang VNPay & \texttt{Payment} ở trạng thái thất bại, không tạo enrollment, frontend hiển thị thông báo thất bại & Đạt \\ \hline
TC-Pay-04 & Callback với chữ ký không hợp lệ (giả mạo) & API từ chối, không cập nhật trạng thái \texttt{Payment} & Đạt \\ \hline
TC-Pay-05 & VNPay gọi lại callback lần thứ hai cho giao dịch đã \texttt{SUCCESS} & Backend trả kết quả thành công ngay, không enroll lại, không ghi đè dữ liệu giao dịch gốc & Đạt \\ \hline
\end{tabular}
\renewcommand{\arraystretch}{1}
\caption{Trường hợp kiểm thử chức năng đăng ký khóa học và thanh toán VNPay}
\label{table:4.4}
\end{table}

Tổng cộng 10 trường hợp kiểm thử thủ công nêu trên đều cho kết quả đạt. Các trường hợp kiểm thử bổ sung cho những chức năng khác (xác thực JWT, enrollment khóa học miễn phí, CRUD nội dung phía quản trị) được em thực hiện theo phương pháp tương tự trong quá trình phát triển, không trình bày chi tiết trong báo cáo.

\section{Triển khai}
\label{section:4.5}

Em triển khai và kiểm thử hệ thống trên môi trường phát triển cục bộ (localhost), gồm ba thành phần chạy song song:

\begin{itemize}
    \item \textbf{Cơ sở dữ liệu MySQL}: chạy qua Docker Compose (\texttt{docker compose up -d mysql}), cổng mặc định \texttt{3306}, dữ liệu được khởi tạo qua migration Flyway và nạp dữ liệu mẫu bằng \texttt{mysql.exe}.
    \item \textbf{Backend Spring Boot}: chạy bằng \texttt{mvn spring-boot:run}, lắng nghe ở cổng \texttt{8080}, cung cấp toàn bộ API dưới tiền tố \texttt{/api/v1}.
    \item \textbf{Frontend React (Vite)}: chạy bằng \texttt{npm run dev}, lắng nghe ở cổng \texttt{5173}, gọi API backend qua biến môi trường \texttt{VITE\_API\_BASE\_URL}.
\end{itemize}

Cấu hình nhạy cảm (chuỗi kết nối CSDL, khóa bí mật JWT, thông tin sandbox VNPay) được em quản lý qua biến môi trường, không hard-code trong mã nguồn.

\end{document}
